\section{讨论与可能的结果}

本章讨论所提出研究的一系列可能结果，及其对学术文献和DeFi协议设计的含义，以及约束结果的解释和推广性的边界条件。与仅呈现单一假设结果的研究不同，本讨论坦承可能呈现多种形态的结果，每种形态均具有独特的含义。

\subsection{情景一：行为过程变量意义重大（H2a获得支持）}

从研究是否对理论有所贡献的角度来看，最为理想的结果是行为过程变量在统计上和经济上均提供了超越常规链上风险指标的显著增量预测能力。若该情景成为现实，研究结果将表明\emph{借贷头寸究竟如何到达其当前状态}，即主动调整的顺序、对风险信号的响应时机以及风险管理行为的一致性，包含了与当前头寸所处状态本身不同但相关的信息，这些信息并未被当前头寸状态所完全捕捉。这一发现将与行为金融经济学中一个不断增长的研究阵营高度一致：该阵营认为投资组合管理过程——而不仅仅是投资组合的最终构成——同样包含关于投资者成熟程度及未来表现的信息。

这一发现的贡献将是三维的。第一，在理论层面，它将提供首次基于头寸粒度级别的系统经验证据，证明前景理论的预测——特别是关于损失厌恶和参照点效应的预测——在DeFi借贷行为中是可观测的。DeFi情境的独特特征——客观的、协议定义的参照点；完全的行为可观测性；不存在传统金融中介所作的混淆因素——将使得这些证据在研究设计角度上特别具有说服力，并对扩展现有行为金融理论在更广泛自然实验语境下的实证计划做出贡献。

第二，在经验层面，它将为链上信用风险工具体系引入一个新的预测变量类别——行为过程变量。虽然这些变量的增量预测能力在绝对量级上可能相对温和（如基于强基准模型进行构建时所常见），但即使是在现有强基准上取得统计显著的边际改善，在DeFi协议背景下也可能具有实际应用意义，因为即使清算预测的边际改善，在考虑涉及金额规模后也将产生实质性的经济影响。

第三，在实践层面，研究结果将提示DeFi协议可以通过行为性早期预警指标来补充其现有纯状态依赖型风险评估机制。例如，协议可以监测特定健康因子区间内借款人主动调整其头寸的比例，并将具有较高风险但异常低调整比例视作即将来临清算压力的组合信号。

\subsection{情景二：行为是系统性的但不具预测性（RQ1获得支持；RQ2不获支持）}

另一种可能的结局是RQ1的行为假说获得支持——借款人确实在头寸接近风险阈值时表现出主动调整的系统性模式，这些模式与前景理论的预测相符合——但RQ2不获支持，原因是行为过程变量未能提供超越常规风险指标的统计显著增量预测能力。在该情景下，研究发现将证明围绕清算阈值的借贷行为是有结构的而非随机的，且该结构可以利用前景理论这一理论框架进行描述，但此类行为结构中包含的信息已经完全被更简单、更直接的风险指标所充分捕捉。健康因子本身，或许补充以近期的变化轨迹和当前市场波动环境，即是包含借贷行为中信用相关信息的充分统计量。

虽然这一情景不及研究更宏大的贡献目标，但其仍将代表一项文献中的宝贵补充。据我们所知，这将是首个系统经验刻画DeFi借贷者如何在头寸风险积累时主动管理——或未能管理——其头寸的研究。即使在缺少预测性成功的情况下，贡献仍在于记录了一组关于借贷行为的风格化事实，这些事实可以为未来理论和经验研究提供启示。

\subsection{情景三：无系统行为亦无预测能力（RQ1和RQ2均不获支持）}

从成果可发表角度最为不利的结果是RQ1和RQ2均不获支持：借款人在头寸接近清算阈值时未表现出主动调整的系统性模式，行为过程变量也无法改善对清算事件的预测。在该情景中，数据将提示DeFi借款人在清算阈值附近的行为要么是随机的，要么完全由价格波动的机械效应和清算机器人活动所驱动，或者由研究无法观测的因素支配。

即使在这一情景下，本研究也同样产出了有价值的零结果。DeFi协议的透明性意味着研究设计和数据是完全可复现的，假说是可证伪的并是事前声明的。然而，在此设定下的零结果，比社会科学中其他经常被用来说明未能复现的因素（例如研究者自由度或发表偏误），更可能反映的实际上是效应的不存在。事实上，对该结果的一种合理解读是，DeFi借贷市场是相对有效的：借款人在清算罚款生成的强财务激励下对上述激励做出了响应，其方式足以由标准风险指标充分近似。这一解读本身便构成了对那些研究DeFi市场效率（在传统金融市场中行为偏差更易得到文献记录）的研究者的兴趣发现。

\subsection{贡献与启示}

本研究所提出的研究有可能从多个层面作出贡献。在方法论层面，它将示范一种如何利用公开可用区块链数据进行详细的、以过程为导向但无需依赖专有数据集或平台合作协议的研究模板。在理论层面，它将在一个新颖的经验设定中提供对行为金融预测的系统性检验，该设定提供了对相关理论构念有着异常干净的识别。在实践层面，它将评估借贷行为是否包含了超越其现有基于状态的评估机制的信用相关信息，从而为DeFi协议的设计提供指导。

本研究的意义超出了特定DeFi借贷范围。本研究试图回答的更广泛问题是，风险管理的\emph{过程}是否包含了超越风险暴露\emph{状态}的信息，这一问题适用于从传统证券市场的保证金交易到公司债务结构管理的广泛金融语境。我们相信，证明这一问题的答案是肯定的，并且可以使用DeFi数据作为概念验证，以为一个终将扩展至其他这些数据远为难以获取的语境的后续研究计划奠定基础。

\subsection{局限性与边界条件}

本研究无论在三种情景的假设实现情况中哪一种发生，均存在若干重要局限性，这些局限性约束了对其结果的解释和推广。

\textbf{用户行为的部分可观测性。} 分析仅限于链上交易，这意味着它无法观测到借款人在中心化交易所或链下的操作。一个在链上表现相对被动的借款人可能正在中心化交易所积极对冲头寸风险，而一个在链上活跃调整头寸的借款人可能出于与该特定头寸风险完全无关的原因如此操作。这种部分可观测性是任何依赖公共区块链数据的固有限制，应当明确承认而非轻描淡写。

\textbf{无法从链上数据推断意图。} 尽管将交易分类为主动借贷行为和被动清算是直接且可复现的，但无法观测到底层操作的具体动机。在健康因子下降后主动追加抵押品的借款人可能是在响应清算的威胁，但也完全可能是在响应与头寸风险相关但概念上不同的其他因素——例如对抵押资产价格未来走势的重新评估，或仅仅是维持既定组合配置的愿望。

\textbf{有限的外部有效性。} 本研究聚焦于一组特定的DeFi借贷协议，并且处于一个特定的市场发展时期。研究发现能在多大程度上推广至其他协议、其他时期乃至其他类型的基于区块链的金融应用仍是未知数，需要进一步重复和扩展方能予以评估。DeFi生态的快速变革——包括协议参数的不断调整、全新借贷机制的引入以及竞争格局的演变——意味着即使在分析时内部有效的研究结论，其适用期限也可能相当有限。

\textbf{不可观测的借款人特征造成的混淆。} 尽管采用包含借款人固定效应的面板数据方法以控制时变不可观测异质性，但仍可能存在随时间变化的不可观测因素——如借款人财务状况、风险偏好或替代流动性来源的变化——会同时影响其对风险积累的行为响应和后续的清算概率，从而产生行为过程变量与清算结果之间的虚假关联。数据的观测性质意味着研究可以识别关联，但无法确立因果关系。

\subsection{伦理考虑}

本研究所提出的研究仅使用不包含个人可识别信息的公开可用链上数据。没有借款人会因为其钱包地址而被联系或重新识别，分析也仅限于在聚合和统计摘要层面进行。然而，本研究触及了关于利用公共区块链数据进行信用评估的更广泛伦理问题。在本研究工作的基础上，协议未来可能会寻求将行为变量纳入风险评估之中，这可能对在其中被涉及的借款人产生不易预期的深远影响。本身值得在此强调的是：主动头寸调整与信用结果之间的关系是在聚合层面进行评估的；若在未仔细留意统计性歧视和公平性问题的情况下，上述估计不能被有意义的解释为个体层面预测。

\subsection{未来研究方向}

本研究拟开辟若干未来研究方向，其中一些可在现阶段预期到，另一些则需待经验结果产生后才能显现。第一，若行为过程变量被证明具有增量预测能力，后续研究可采用更细致的分解分析或结构估计方法研究驱动此一关系的机制。第二，本研究开发的框架可推广至运行于不同区块链网络上的其他DeFi借贷协议，从而提供关于研究结果能否在不同协议设计和市场环境中进行推广的证据。第三，未来研究可探讨将行为过程变量纳入协议风险参数对福利的影响，既考虑改进风险评估的潜在收益，也考虑协议复杂性提高带来的潜在成本。第四，未来研究可探讨在DeFi中表现出系统性主动头寸管理模式（的借款人），条件允许的前提下在传统信用市场中是否也表现出不同性质的内生（信贷）体现。
